%\makeglossaries
\makenoidxglossaries

% Danh mục thuật ngữ và từ viết tắt dùng trong đồ án
\newglossaryentry{mac}{
    type=\acronymtype,
    name={MAC},
    description={Phép nhân--cộng tích lũy (Multiply--Accumulate)},
    first={nhân--cộng tích lũy (Multiply--Accumulate - MAC)}
}

\newglossaryentry{simd}{
    type=\acronymtype,
    name={SIMD},
    description={Mô hình một lệnh xử lý nhiều dữ liệu (Single Instruction, Multiple Data)},
    first={Single Instruction, Multiple Data (SIMD)}
}

\newglossaryentry{dsp}{
    type=\acronymtype,
    name={DSP},
    description={Khối hoặc kỹ thuật xử lý tín hiệu số (Digital Signal Processing/Processor)},
    first={Digital Signal Processing/Processor (DSP)}
}

\newglossaryentry{rvv}{
    type=\acronymtype,
    name={RVV},
    description={Mở rộng véc-tơ của RISC-V (RISC-V Vector Extension)},
    first={RISC-V Vector Extension (RVV)}
}

\newglossaryentry{soc}{
    type=\acronymtype,
    name={SoC},
    description={Hệ thống trên chip (System-on-Chip)},
    first={hệ thống trên chip (System-on-Chip - SoC)}
}

\newglossaryentry{tflm}{
    type=\acronymtype,
    name={TensorFlow Lite Micro},
    description={Khung suy luận học máy dành cho vi điều khiển và hệ nhúng tài nguyên hạn chế},
    first={TensorFlow Lite Micro}
}

\newglossaryentry{litert}{
    type=\acronymtype,
    name={LiteRT},
    description={Runtime suy luận trên thiết bị biên của Google AI Edge; trong báo cáo được dùng làm nguồn đặc tả lượng tử hóa 8 bit},
    first={LiteRT}
}

\newglossaryentry{mnist}{
    type=\acronymtype,
    name={MNIST},
    description={Bộ dữ liệu chữ số viết tay thường dùng để đánh giá mô hình phân loại ảnh đơn giản},
    first={MNIST}
}

\newglossaryentry{fc}{
    type=\acronymtype,
    name={FC},
    description={Lớp kết nối đầy đủ trong mạng nơ-ron (Fully Connected)},
    first={fully connected (FC)}
}

\newglossaryentry{relu}{
    type=\acronymtype,
    name={ReLU},
    description={Hàm kích hoạt tuyến tính chỉnh lưu (Rectified Linear Unit)},
    first={Rectified Linear Unit (ReLU)}
}

\newglossaryentry{tpu}{
    type=\acronymtype,
    name={TPU},
    description={Bộ xử lý chuyên dụng cho tensor do Google giới thiệu (Tensor Processing Unit)},
    first={Tensor Processing Unit (TPU)}
}

\newglossaryentry{tops}{
    type=\acronymtype,
    name={TOPS},
    description={Đơn vị đo thông lượng, nghìn tỷ phép toán mỗi giây (Tera Operations Per Second)},
    first={Tera Operations Per Second (TOPS)}
}
